% !TeX root = ../../Book2.tex
\section{多线性限制的几何与尺度方法}
\label{b2:sec:2505}

前面的横截模型解释了不同方向怎样改善乘积估计。把这种增益用于许多帽和许多空间区域，还需要尺度递推、波包分解以及对几何组合的控制。本节先证明其中可直接计算的步骤，再陈述抛物面解耦这一深层输入。递推式的证明、单帽展开和多帽组合是不同的任务，应分别检查它们需要的条件。

\subsection{尺度归纳}
\label{b2:subsec:250501}

设 \(A(R)\) 表示半径 \(R\) 球上的某个规范化最佳常数；若定理允许任意球心或一族曲面，便应先对这些参数取上确界。尺度归纳希望把 \(A(R)\) 与较小尺度的常数联系起来。一种常见模型是
\[
 A(R)\leq C_\eta R^\eta A(R^{1/2}).
\]
右侧不只有幂次损失：每迭代一次，还会多出一个 \(C_\eta\)。下面把两者都算出。

\Needspace{6\baselineskip}
\begin{lemma}[平方根尺度递推]\label{b2:lem:25-scale-iteration}
设 \(A:[1,\infty)\to[0,\infty)\) 为有限值函数，满足
\[
 M\coloneqq\sup_{1\leq R\leq2}A(R)<\infty.
\]
假设对每个 \(\eta>0\)，存在与 \(R\) 无关的有限常数 \(C_\eta\geq1\)，使
\[
 A(R)\leq C_\eta R^\eta A(R^{1/2})\quad(R>2).
\]
则对每个 \(\varepsilon>0\)，存在有限常数 \(K_\varepsilon\)，使
\[
 A(R)\leq K_\varepsilon R^\varepsilon\quad(R\geq1).
\]
常数 \(K_\varepsilon\) 可以依赖 \(M\) 及相应的 \(C_\eta\)，不要求它在 \(\varepsilon\downarrow0\) 时保持有界。
\end{lemma}
\begin{proof}
固定 \(\eta>0\) 及 \(R>2\)，令 \(m\) 为满足 \(R^{2^{-m}}\leq2\) 的最小正整数。前 \(m\) 步均可应用递推，故
\[
 \begin{aligned}
 A(R)
 &\leq C_\eta^m
       R^{\eta(1+2^{-1}+\cdots+2^{-(m-1)})}
       A(R^{2^{-m}})\\
 &\leq M C_\eta^m R^{2\eta}.
 \end{aligned}
\]
另一方面，
\[
 m\leq1+\log_2\!\left(\frac{\log R}{\log2}\right).
\]
记 \(b_\eta=\log_2 C_\eta\geq0\)，则
\[
 C_\eta^m
 \leq C_\eta\left(\frac{\log R}{\log2}\right)^{b_\eta}.
\]
所以固定 \(\eta\) 后，全部常数累积只是一个固定的对数幂。

现在给定 \(\varepsilon>0\)，先取 \(\eta=\varepsilon/4\)。这样 \(C_\eta\) 和 \(b_\eta\) 已经固定。函数
\((\log R)^{b_\eta}R^{-\varepsilon/2}\)
在 \(R\geq2\) 上有界：令 \(t=\log R\)，它就是 \(t^{b_\eta}e^{-\varepsilon t/2}\)，指数衰减控制固定幂次。因此
\[
 A(R)\leq K_\varepsilon
          R^{2\eta+\varepsilon/2}
       =K_\varepsilon R^\varepsilon.
\]
初始区间 \(1\leq R\leq2\) 由 \(M\) 控制，增大常数即可。
\end{proof}

顺序在这里很重要：先选 \(\eta\)，再吸收由 \(C_\eta\) 决定的对数幂。若让 \(\eta\) 随 \(R\) 改变，却不知道 \(C_\eta\) 怎样增长，上述吸收便没有统一常数。这个引理也不负责证明具体限制问题所需的递推式；它只说明一旦递推和初始尺度界成立，怎样正确结束迭代。

\ProseParagraphBreak
每一步还必须留在同一允许的曲面类中，包括参数化的导数界、支集大小和横截下界。\cref{b2:prop:25-parabolic-rescaling}给出抛物面的精确重标度；一般曲面的小片重标度则要控制偏离二次模型的误差。如果归纳后曲率或横截常数不断恶化，就不能把下一步仍写成同一个 \(A(R^{1/2})\)，更不能只凭幂次几何级数断言得到任意小损失。

\subsection{波包}
\label{b2:subsec:250502}

在空间半径 \(R\) 上，一个宽度 \(R^{-1/2}\) 的频率帽对应宽约 \(R^{1/2}\)、长约 \(R\) 的管。这与\cref{b2:fig:23-cap-tube}及\cref{b2:prop:23-parabolic-packet}的尺度一致，但单个光滑帽函数还不代表一般 \(L^2\) 输入。下面用帽内 Fourier 级数把一般输入拆成位置不同的波包，并保留系数的平方能量。

\ProseParagraphBreak
为使归一化可直接检查，这一构造使用抛物面的参数测度；面积测度的光滑密度在固定参数片上可以吸收到输入中，但相应的范数比较应随之注明。

\Needspace{6\baselineskip}
\begin{proposition}[单帽的波包展开]\label{b2:prop:25-single-cap-packets}
设整数 \(n\geq1\)，\(R\geq1\)、\(\delta=R^{-1/2}\)，且 \(\xi_0\in\mathbb R^n\)、\(|\xi_0|\leq1\)。取
\[
 g\in L^2(\mathbb R^n),\quad
 \operatorname{supp}g\subset
       \xi_0+\delta[-1/2,1/2]^n,
\]
并以参数测度定义
\[
 E_{\mathrm{par}}g(x',t)
 =\int_{\mathbb R^n}
       e^{2\pi\mathrm i(x'\cdot\xi+t|\xi|^2)}
       \cdot g(\xi)\,\mathrm d\xi.
\]
固定 \(\rho\in C_c^\infty((-1,1)^n)\)，在 \([-1/2,1/2]^n\) 上恒为一。则存在系数 \(c_k\)、\(k\in\mathbb Z^n\)，及只依赖 \(\xi_0,\delta,\rho,k\) 的频率函数 \(\psi_k\)，使
\[
 g=\sum_{k\in\mathbb Z^n}c_k\psi_k
       \quad\text{于 }L^2(\mathbb R^n),\quad
 \sum_k|c_k|^2=\|g\|_2^2.
\]
若 \(W_k=E_{\mathrm{par}}\psi_k\)，则按有限立方格点部分和取极限时，
\[
 E_{\mathrm{par}}g=\sum_kc_kW_k
\]
在整个 \(\mathbb R^{n+1}\) 上一致收敛。对每个 \(N>0\)，在 \(|t|\leq R\) 内还有
\[
 |W_k(x',t)|
 \leq C_{N,n,\rho}\delta^{n/2}
  \left(1+\left|\delta(x'+2t\xi_0)+k/2\right|\right)^{-N}.
\]
可选一个只依赖维数的固定 \(A_0\)，使对任意 \(L>0\)，
\[
 \sup_{|(x',t)|\leq R}
 \left|E_{\mathrm{par}}g(x',t)
       -\sum_{|k|\leq A_0\sqrt R}c_kW_k(x',t)\right|
 \leq C_{L,n,\rho}R^{-L}\|g\|_2.
\]
所有常数均与 \(R,\xi_0,g\) 无关。
\end{proposition}
\begin{proof}
将帽拉回单位盒，令 \(h(u)=g(\xi_0+\delta u)\)。它在
\(Q=[-1,1]^n\) 上支于较小的盒 \([-1/2,1/2]^n\)。正交基
\[
 e_k(u)=2^{-n/2}e^{\pi\mathrm i k\cdot u},\quad k\in\mathbb Z^n,
\]
给出 Fourier 级数
\[
 h=\sum_ka_ke_k\quad\text{于 }L^2(Q),\quad
 \sum_k|a_k|^2=\|h\|_{L^2(Q)}^2
             =\delta^{-n}\|g\|_2^2.
\]
令
\[
 c_k=\delta^{n/2}a_k,\quad
 \psi_k(\xi)=2^{-n/2}\delta^{-n/2}
       \rho\!\left(\frac{\xi-\xi_0}{\delta}\right)
       \cdot e^{\pi\mathrm i k\cdot(\xi-\xi_0)/\delta}.
\]
由于 \(\rho\cdot h=h\)，把 Fourier 级数乘以有界函数 \(\rho\)，再作换元，就得到所述 \(L^2\) 重构与精确系数能量。每个频率函数支于放大一倍的帽，且
\[
 \|\psi_k\|_2=2^{-n/2}\|\rho\|_2.
\]
它们有固定的 \(L^2\) 大小，但并不因这个归一化就彼此正交；系数能量等式来自乘 \(\rho\) 之前的正交 Fourier 基。

换元 \(\xi=\xi_0+\delta u\) 给出原波包的具体公式
\[
 \begin{aligned}
 W_k(x',t)
 ={}&2^{-n/2}\delta^{n/2}
       e^{2\pi\mathrm i(x'\cdot\xi_0+t|\xi_0|^2)}\\
 &{}\cdot
 \int_{\mathbb R^n}\rho(u)\cdot
       e^{2\pi\mathrm i[
          \delta(x'+2t\xi_0)\cdot u
          +t\delta^2|u|^2+k\cdot u/2]}\,\mathrm du.
 \end{aligned}
\]
频率部分和的误差支于固定的放大帽，其 \(L^2\) 范数趋零，故其 \(L^1\) 范数也趋零。延拓积分核的模为一，因此输出误差的一致范数至多为这个 \(L^1\) 误差。这证明全空间的一致收敛，并没有假定一般 \(L^2\) 函数的 Fourier 级数逐点收敛。

若 \(|t|\leq R\)，则 \(|t\delta^2|\leq1\)。振幅
\(\rho(u)\cdot e^{2\pi\mathrm i t\delta^2|u|^2}\)
的支集固定，各固定阶导数一致有界。对剩余的线性频率
\[
 z+k/2,\quad z=\delta(x'+2t\xi_0),
\]
反复分部积分，即得所述任意阶衰减。

还需把这些尾部实际求和。对任意 \(z\in\mathbb R^n\)，格点间距固定，半径 \(2^j\) 的球内至多有 \(C_n2^{jn}\) 个点。因此当 \(2N>n\) 时，
\[
 \sup_z\sum_{k\in\mathbb Z^n}
       (1+|z+k/2|)^{-2N}<\infty.
\]
离散 Cauchy--Schwarz 与系数能量给出
\[
 \sum_k|c_k|\cdot|W_k(x',t)|
 \leq C_{N,n,\rho}\delta^{n/2}\|g\|_2
 \quad(|t|\leq R).
\]
所以该时间片内的波包和还逐点绝对收敛。

在球 \(|(x',t)|\leq R\) 内，因 \(|\xi_0|\leq1\)，有
\(|z|\leq3\sqrt R\)。取固定 \(A_0\geq12\)，若 \(|k|>A_0\sqrt R\)，则
\[
 |z+k/2|\geq|k|/4.
\]
同样按格点环带求和，得到
\[
 \begin{aligned}
 \sum_{|k|>A_0\sqrt R}|c_k|\cdot|W_k(x',t)|
 &\leq C_{N,n,\rho}\delta^{n/2}\|g\|_2
        R^{n/4-N/2}\\
 &=C_{N,n,\rho}R^{-N/2}\|g\|_2.
 \end{aligned}
\]
取 \(N>\max\{n/2,2L\}\)，便得到命题的局部误差界。
\end{proof}

在这个构造中，主要管区域为
\[
 |x'+2t\xi_0+k/(2\delta)|\lesssim\delta^{-1},\quad
 |t|\leq R.
\]
其宽度约为 \(R^{1/2}\)，长度约为 \(R\)，长轴沿抛物面的法向 \((-2\xi_0,1)\)。帽的中心决定方向，格点 \(k\) 决定位置；同一帽中的管彼此平行。空间球内只需保留 \(O(R^{n/2})\) 个位置指标，其余项可以按任意预先指定的负幂控制。

\ProseParagraphBreak
这里有三个不同的结论。频率重构在 \(L^2\) 中成立；紧频率支集又使延拓部分和在全空间一致收敛；管外快速衰减及上述有限格点截断则限于 \(|t|\leq R\) 或 \(|(x',t)|\leq R\)。最后一项不能无条件推广到任意大时间。波包并非管的示性函数：公式中的相位与系数仍须保留，截到粗管时产生的尾部也须按需要的阶数求和。

\ProseParagraphBreak
这只是单张帽的展开。若把许多帽相加，各帽之间有不同的法向，不能从一张帽的系数能量等式推出 BCT 或解耦。为了估计这些组合，接下来需要同时看管与空间分区的关系，以及不同帽输出的范数怎样相加。

\subsection{多项式分割}
\label{b2:subsec:250503}

多项式分割选择一个多项式，把大部分空间分成若干胞腔，并将零集的邻域单独保留。它希望让各胞腔承担可控的质量，同时利用管的穿越次数限制重复计算。选择分割多项式及控制各胞腔质量属于额外的存在性结论，不能由直线的根数直接推出。

\ProseParagraphBreak
\cref{b2:lem:23-polynomial-line-crossing}证明：次数至多 \(D\) 的非零多项式，其零集之外的胞腔，至多有 \(D+1\) 个能被同一条不全含于零集的直线经过。对于半径为 \(\rho\) 的粗管，要把胞腔缩到与零集距离大于 \(\rho\) 的部分，才能把管的相遇转回中心直线的穿越。原因是管中这样的点到中心线的连接线段避开零集，故其中心线点仍在同一胞腔。未经缩减的胞腔与零集邻域，不由这个次数界控制。

即使多项式只有一次，零集邻域中的管也会表现出两种不同情形。平面模型足以把角度和停留长度算清。

\Needspace{6\baselineskip}
\begin{proposition}[管与平面厚邻域]\label{b2:prop:25-plane-wall}
设整数 \(d\geq2\)，\(\nu,v\in S^{d-1}\)，\(a\in\mathbb R^d\)、\(b\in\mathbb R\)，且 \(L,r,\rho>0\)。令
\[
 H=\{\,x:x\cdot\nu=b\,\},\quad
 W_r(H)=\{\,x:|x\cdot\nu-b|\leq r\,\},
\]
\[
 \gamma(s)=a+sv\quad(0\leq s\leq L),\quad
 T_\rho=\bigcup_{0\leq s\leq L}\overline{B(\gamma(s),\rho)}.
\]
记可能与平面厚邻域相遇的轴参数集合为
\[
 I=\{\,s\in[0,L]:
       \overline{B(\gamma(s),\rho)}\cap W_r(H)\neq\varnothing\,\}.
\]
以一维 Lebesgue 测度记 \(|I|\)。若 \(v\cdot\nu\neq0\)，则
\[
 |I|\leq\min\left\{L,\frac{2(r+\rho)}{|v\cdot\nu|}\right\}.
\]
若 \(v\cdot\nu=0\)，则 \(I=[0,L]\) 或 \(I=\varnothing\)，前一种情形恰好对应 \(|a\cdot\nu-b|\leq r+\rho\)。特别地，若轴线在 \(H\) 内且 \(\rho\leq r\)，则整条管 \(T_\rho\) 都包含于 \(W_r(H)\)。
\end{proposition}
\begin{proof}
球 \(\overline{B(\gamma(s),\rho)}\) 与平面厚邻域相遇，当且仅当其球心到平面的距离至多为 \(r+\rho\)。于是
\[
 I=\{\,s\in[0,L]:
       |a\cdot\nu-b+s(v\cdot\nu)|\leq r+\rho\,\}.
\]
若斜率 \(v\cdot\nu\neq0\)，不等式在整个实轴上给出长度
\(2(r+\rho)/|v\cdot\nu|\) 的区间，与 \([0,L]\) 相交后得到所述上界。若斜率为零，条件与 \(s\) 无关，所以参数集合全有或全无。轴线位于平面内时，管中每点到该平面的距离不超过 \(\rho\)，最后一项随即成立。
\end{proof}

若 \(\vartheta\) 是管轴与平面的夹角，则 \(|v\cdot\nu|=|\sin\vartheta|\)。取波包尺度 \(L=R\)、\(r\) 与 \(\rho\) 都约为 \(R^{1/2}\)，固定正角度给出的停留长度只有 \(O(R^{1/2})\)；当角度小到 \(R^{-1/2}\) 的量级时，上界可以达到整条管的长度 \(R\)。平行而贴墙的管更可能一直留在墙内，不能把它作为普通横穿管处理。

\ProseParagraphBreak
这说明零集邻域中的分析不能只问“中心线与零集交了几次”。横穿部分可利用角度下界缩短停留区域；近切向部分则可能沿低维几何聚集，需要不同的估计。对高次零集，局部切平面会随位置改变，还会有奇点及多个局部片，平面模型本身不处理这些困难。

\ProseParagraphBreak
Guth 的《A Restriction Estimate Using Polynomial Partitioning》\cite{Guth2016Restriction}给出了这种方法的正式限制应用。多项式次数一方面增加胞腔数，另一方面也增加穿越预算和零集邻域的复杂度，必须放进尺度递推中选择。本节没有重建分割定理、切向管分析或完整限制范围证明，因而不从上述次数与角度计算推出新的线性限制估计。

\subsection{解耦}
\label{b2:subsec:250504}

单帽波包保留帽内的相干性；解耦则控制许多帽叠加后的输出。下面采用\term[jieou-dingli]{解耦定理}[decoupling theorem]的抛物面临界版本。帽的尺度、分割函数和空间权都属于陈述的一部分。

\Needspace{6\baselineskip}
\begin{theorem}[深层结构定理：抛物面临界解耦]\label{b2:thm:25-decoupling}
设整数 \(d\geq2\)、\(R\geq1\)，令
\[
 S=\{\,(\xi',|\xi'|^2):|\xi'|\leq1\,\}\subset\mathbb R^d,\quad
 p_c=\frac{2(d+1)}{d-1},\quad \delta=R^{-1/2}.
\]
曲面取面积测度 \(\sigma\)，并定义
\[
 Eg(x)=\int_Se^{2\pi\mathrm i x\cdot\xi}
       \cdot g(\xi)\,\mathrm d\sigma(\xi)
 \quad(g\in L^2(S,\sigma)).
\]
在参数域上取尺度 \(\delta\) 的规则方格帽覆盖，并固定非负光滑单位分解 \((\chi_\theta)_\theta\)：每个 \(\chi_\theta\) 支于相应方格的固定倍数放大，其支集重叠次数一致有界，
\[
 \sum_\theta\chi_\theta=1\quad\text{于 }S,\quad
 |\partial^\alpha\chi_\theta|\leq C_\alpha\delta^{-|\alpha|}.
\]
导数在参数坐标中计算，\(C_\alpha\) 与 \(R,\theta\) 无关。记
\[
 E_\theta g\coloneqq E(\chi_\theta\cdot g).
\]
若 \(B_R=B(c,R)\subset\mathbb R^d\)，令
\[
 w_{B_R}(x)=\left(1+\frac{|x-c|}{R}\right)^{-10d},\quad
 \|F\|_{L^{p_c}(w_{B_R})}
 =\left(\int_{\mathbb R^d}|F(x)|^{p_c}
          w_{B_R}(x)\,\mathrm dx\right)^{1/p_c}.
\]
则对每个 \(\varepsilon>0\)，存在 \(C_\varepsilon<\infty\)，使
\[
 \|Eg\|_{L^{p_c}(B_R)}
 \leq C_\varepsilon R^\varepsilon
       \left(\sum_\theta
          \|E_\theta g\|_{L^{p_c}(w_{B_R})}^2
       \right)^{1/2}.
\]
常数可以依赖 \(d,\varepsilon\) 及固定分割的重叠、导数界，与 \(R,c,g\) 无关。
\end{theorem}
\begin{proofsketch}
频率薄邻域中的全空间解耦经空间局部化得到上述延拓形式；固定的光滑帽分割用于定义各项，有限重叠只改变常数。完整证明结合多线性限制、抛物重标度和反复尺度迭代，最后保留任意小幂次损失。这里略去具体的多尺度递推与闭合，不能仅凭\cref{b2:lem:25-scale-iteration}替代那些步骤。

正式来源为 Bourgain 与 Demeter 的《The Proof of the \(\ell^2\) Decoupling Conjecture》%
\cite[Theorem 1.1; equation (7), pp.~352--355]{BourgainDemeter2015Proof}。
该版本要求正定第二基本形式，临界指数如上。标准抛物面满足这些条件；不能把同一结论未经修改套到第二基本形式不定或退化的曲面。
\end{proofsketch}

定义中的 \(E_\theta g\) 始终使用 \(\chi_\theta\cdot g\)，所以
\(Eg=\sum_\theta E_\theta g\) 精确成立。若改用互不相交的可测帽示性函数，应从头采用那种分割的版本；不能既把帽当作重叠覆盖，又让它们的示性函数相加充当一。权 \(w_{B_R}\) 在球内有统一正下界，在球外保留快速衰减尾部，这与波包没有严格管支集相适应。换成双边可比的权只改变常数，不能任意删去球外部分而保持同一陈述。

\ProseParagraphBreak
还应区分两个容易混写的右端。对有限个函数 \(F_\theta\) 及任意正测度 \(\mu\)，当 \(2\leq p<\infty\) 且各项属于 \(L^p(\mu)\) 时，Minkowski 不等式给出
\[
 \begin{aligned}
 \left\|\left(\sum_\theta|F_\theta|^2\right)^{1/2}\right\|_{L^p(\mu)}^2
 &=\left\|\sum_\theta|F_\theta|^2\right\|_{L^{p/2}(\mu)}\\
 &\leq\sum_\theta
       \bigl\||F_\theta|^2\bigr\|_{L^{p/2}(\mu)}
 =\sum_\theta\|F_\theta\|_{L^p(\mu)}^2.
 \end{aligned}
\]
因此点态平方函数的范数至多为各函数范数的 \(\ell^2\) 和。把解耦定理右端换成前者，会提出一个通常更强的估计，不能由显示的解耦不等式直接推出。

\ProseParagraphBreak
一个有限模型可把差别算出。取 \(\mathbb R^d\) 中 \(N\) 个互不相交、体积为一的方块 \(Q_1,\ldots,Q_N\)，令 \(F_j=\mathbf1_{Q_j}\)。则
\[
 \left\|\left(\sum_{j=1}^N|F_j|^2\right)^{1/2}\right\|_p=N^{1/p},
 \quad
 \left(\sum_{j=1}^N\|F_j\|_p^2\right)^{1/2}=N^{1/2}.
\]
在临界 \(p=p_c\) 处，两者的比值为 \(N^{1/(d+1)}\)。这个例子只比较两种范数结构，没有把方块示性函数宣称为曲面帽的延拓。若各个 \(F_j\) 完全相同，两者又相等；差别取决于各项在空间中的分布，不能用一句“正交”概括。

\ProseParagraphBreak
在 \(p=2\) 时，上述两种范数由积分的线性性相等；非平方指数处则存在实际差距。抛物面临界解耦的内容是以 \(R^\varepsilon\) 控制许多频率方向的叠加，并保留加权球和指定帽尺度。单帽的系数能量、平面穿越计算及尺度递推引理各解决其中一部分问题，不能单独代替这项多帽估计。
